\documentclass[a4paper]{article}

% Basic packages
\usepackage[fontset=fandol]{ctex}
\usepackage{amsmath, amssymb}
\usepackage{graphicx}
\usepackage[margin=2.5cm]{geometry}
\usepackage[most]{tcolorbox}
\usepackage{etoolbox}
\usepackage{listings}
\usepackage{hyperref}
\usepackage{booktabs}
\usepackage{subcaption}
\usepackage{float}
\usepackage{tikz}

% ===== Highlight boxes =====
\newtcolorbox{knowledgebox}[1]{
    enhanced, colback=blue!5!white, colframe=blue!75!black, colbacktitle=blue!75!black,
    coltitle=white, fonttitle=\bfseries, title={#1}, attach boxed title to top left={yshift=-2mm, xshift=2mm},
    boxrule=1pt, sharp corners
}

\newtcolorbox{importantbox}[1]{
    enhanced, colback=yellow!10!white, colframe=yellow!80!black, colbacktitle=yellow!80!black,
    coltitle=black, fonttitle=\bfseries, title={#1}, sharp corners
}

\newtcolorbox{warningbox}[1]{
    enhanced, colback=red!5!white, colframe=red!75!black, colbacktitle=red!75!black,
    coltitle=white, fonttitle=\bfseries, title={#1}, sharp corners
}

% ===== Code listings =====
\lstset{
    basicstyle=\ttfamily\small,
    keywordstyle=\color{blue},
    stringstyle=\color{red!60!black},
    commentstyle=\color{green!60!black},
    breaklines=true,
    frame=single,
    numbers=left,
    numberstyle=\tiny\color{gray},
    captionpos=b,
    extendedchars=false
}

% ===== Front-page metadata =====
\newcommand{\notetitle}{损失函数也会影响优化难度}
\newcommand{\notesubtitle}{李宏毅《机器学习》2025 · 类神经网络训练不起来怎么办（四）}
\newcommand{\noteauthors}{基于李宏毅老师课程整理}
\newcommand{\notedate}{\today}
\newcommand{\videochannel}{李宏毅深度学习课堂（Bilibili）}
\newcommand{\videoduration}{约 20 分钟}
\newcommand{\videourl}{https://www.bilibili.com/video/BV1TAtwzTE1S?p=7}

\begin{document}

\pagestyle{empty}

% ===== Cover page =====
\begin{center}
    \vspace*{1.5cm}
    {\Huge\bfseries \notetitle\par}
    \vspace{0.4cm}
    {\LARGE \notesubtitle\par}
    \vspace{0.8cm}
    \includegraphics[width=0.62\textwidth]{video.jpg}\par
    \vspace{0.8cm}
    {\large \noteauthors\par}
    {\large \notedate\par}
    \vspace{0.8cm}
    \begin{tcolorbox}[width=0.8\textwidth,center,sharp corners,
                      colback=black!3!white,colframe=black!50!white]
        \centering
        \textbf{视频信息}\\[3pt]
        频道：\videochannel \\
        时长：\videoduration \\
        链接：\href{\videourl}{\videourl}
    \end{tcolorbox}
\end{center}

\newpage

% ===== TOC =====
\tableofcontents
\newpage

\pagestyle{plain}

% =====================================================================
\section{引言：分类问题能不能直接当回归做？}
% =====================================================================

这一讲是分类问题的短版说明。前几讲已经建立了回归模型、梯度下降、批次、动量、自适应学习率等观念；本讲把焦点换成一个很实际的问题：\textbf{做分类时，输出、softmax 和 loss function 应该怎么搭配？}

最直接的想法是把分类也当成回归：输入 $x$，模型输出一个数值 $y$，然后把类别也编号成数字，例如 $1=\text{class 1}$、$2=\text{class 2}$、$3=\text{class 3}$，再让 $y$ 尽量接近正确类别编号。

\begin{figure}[H]
    \centering
    \includegraphics[width=0.86\textwidth]{figures/fig01_class_as_regression.jpg}
    \caption{把分类当回归做的朴素想法：模型仍然输出单一标量 $y$，再让 $y$ 接近类别编号。但这等于把 class 1、class 2、class 3 放在一条有距离关系的数轴上。\protect\footnotemark}
    \label{fig:class_as_regression}
\end{figure}
\footnotetext{视频画面与字幕时间区间：01:25--02:08。}

\begin{warningbox}{类别编号会偷偷加入``相似度''假设}
    如果把 class 1 编成 $1$、class 2 编成 $2$、class 3 编成 $3$，模型会隐含地认为 class 1 比较接近 class 2，而比较远离 class 3。这个假设在某些有序类别中可能合理，例如年级、等级、评分；但在一般分类中通常不成立。例如猫、狗、车之间没有天然的一维顺序，把它们硬塞进 $1,2,3$ 会给模型错误的几何结构。
\end{warningbox}

这里的重点不是说``类别永远不能编号''，而是要看编号是否承载了真实语义。若任务本身是有序预测，例如一年级、二年级、三年级，编号有可能反映真实顺序；但若任务是手写数字、宝可梦种类、动物类别，类别编号往往只是资料集作者随手排的 ID。模型若把这个 ID 当成连续数值，就会学习到不存在的顺序关系。

因此，分类问题通常不用一个数字表示类别，而是把类别表示为 \textbf{one-hot vector}。

% =====================================================================
\section{One-hot 标签与多输出网络}
% =====================================================================

若共有三个类别，则正确标签 $\hat{\boldsymbol{y}}$ 可以写成三维 one-hot vector：

\[
    \text{class 1}: \begin{bmatrix}1\\0\\0\end{bmatrix},\qquad
    \text{class 2}: \begin{bmatrix}0\\1\\0\end{bmatrix},\qquad
    \text{class 3}: \begin{bmatrix}0\\0\\1\end{bmatrix}.
\]

这种表示方法的好处是：每一类都占一个独立维度，类别之间不再因为编号大小产生人为远近。若用欧氏距离看，三个 one-hot 向量两两之间的距离相同，比较符合一般分类任务中``类别只是不同，不一定有顺序''的直觉。

\begin{knowledgebox}{One-hot 做了什么，又没有做什么}
    One-hot 的作用是\textbf{把类别 ID 从连续数轴上解放出来}。它告诉模型：这些类别是彼此不同的离散选项，而不是 $1<2<3$ 的连续量。但 one-hot 本身并不知道类别之间的语义相似性。例如 class 1 与 class 2 是否真的相近，one-hot 不表达；它只负责给 supervised learning 一个干净的目标格式。若未来想表达类别语义，那是 embedding、metric learning 或 label hierarchy 的问题，不是这讲要解决的分类基本设定。
\end{knowledgebox}

\begin{figure}[H]
    \centering
    \includegraphics[width=0.86\textwidth]{figures/fig02_one_hot_network.jpg}
    \caption{当标签变成三维 one-hot vector，网络也要输出三个数值 $y_1,y_2,y_3$。做法并不神秘：把原本输出一个数值的最后一层改成输出多个数值，每个输出都有自己的一组权重和偏置。\protect\footnotemark}
    \label{fig:one_hot_network}
\end{figure}
\footnotetext{视频画面与字幕时间区间：04:55--05:26。}

以一个隐藏层网络为例，回归任务可以写成

\[
    y = b + \boldsymbol{c}^{T}\sigma(\boldsymbol{b}+W\boldsymbol{x}),
\]

其中 $y$ 是一个标量。分类时，最后一层输出不再是单一数字，而是向量：

\[
    \boldsymbol{y}
    = \boldsymbol{b}' + W'\sigma(\boldsymbol{b}+W\boldsymbol{x})
    =
    \begin{bmatrix}
    y_1\\y_2\\y_3
    \end{bmatrix}.
\]

这些 $y_i$ 还不是概率，也不一定介于 $0$ 与 $1$ 之间。它们只是网络最后一层吐出的原始分数，常称为 \textbf{logits}。下一步，就是用 softmax 把 logits 变成可与 one-hot 标签比较的输出。

\begin{knowledgebox}{从标量输出到向量输出}
    分类网络的关键改变不是``网络变成另一种东西''，而是\textbf{最后一层的输出维度变成类别数}。若有 $K$ 个类别，网络输出 $K$ 个 logits；第 $i$ 个 logit 可以理解成模型对第 $i$ 类的未归一化支持程度。
\end{knowledgebox}

这个改动也解释了为什么分类模型的最后一层常写成一个矩阵乘法。若隐藏层表示为 $\boldsymbol{h}=\sigma(\boldsymbol{b}+W\boldsymbol{x})$，最后一层就是

\[
    \boldsymbol{y}=W'\boldsymbol{h}+\boldsymbol{b}',
\]

其中 $W'$ 的每一行对应一个类别的打分器。换句话说，分类不是只训练一个回归器，而是同时训练 $K$ 个类别打分器，再把这些分数交给 softmax 统一比较。

\begin{importantbox}{分类模型里的四个对象}
    一条标准多分类链路里，最好把下面四个对象分清楚：
    \begin{itemize}
        \item $\boldsymbol{x}$：输入特征，例如影像、语音特征、文字向量等。
        \item $\boldsymbol{y}$：网络输出的 logits，尚未归一化，可以是任意实数。
        \item $\boldsymbol{y}'$ 或 $\boldsymbol{p}$：softmax 后的输出，所有分量为正且总和为 $1$。
        \item $\hat{\boldsymbol{y}}$ 或 $\boldsymbol{q}$：正确答案的 one-hot 标签。
    \end{itemize}
    训练时真正比较的是 $\boldsymbol{p}$ 与 $\boldsymbol{q}$；但优化时真正被更新的是产生 logits $\boldsymbol{y}$ 的网络参数。
\end{importantbox}

% =====================================================================
\section{Softmax：把 logits 变成可比较的分布}
% =====================================================================

one-hot 标签里的元素只会是 $0$ 或 $1$，而 logits 可以是任意实数。为了让网络输出与标签在同一种尺度上比较，分类模型通常在 logits 后接一个 \textbf{softmax}：

\begin{figure}[H]
    \centering
    \includegraphics[width=0.86\textwidth]{figures/fig03_softmax_position.jpg}
    \caption{分类模型的输出链路：网络先产生 logits $\boldsymbol{y}$，再经过 softmax 得到 $\boldsymbol{y}'$，最后计算 $\boldsymbol{y}'$ 与 one-hot 标签 $\hat{\boldsymbol{y}}$ 之间的误差。\protect\footnotemark}
    \label{fig:softmax_position}
\end{figure}
\footnotetext{视频画面与字幕时间区间：06:56--07:40。}

softmax 的定义为

\[
    y'_i = \frac{\exp(y_i)}{\sum_j \exp(y_j)}.
\]

它做了两件重要的事：

\begin{itemize}
    \item 先用 exponential 把每个 logit 变成正数，因此 $y'_i>0$。
    \item 再除以所有 exponential 的和，因此 $\sum_i y'_i = 1$。
\end{itemize}

所以 $\boldsymbol{y}'$ 可以被理解为一个归一化后的类别分布。严格地说，它不是模型``天生知道的真实概率''，而是我们设计出来、便于训练和解释的输出形式。

Softmax 还有一个很重要的性质：\textbf{它只在乎 logits 的相对差距，不在乎整体平移}。若对所有 logits 同时加上同一个常数 $c$，

\[
    \frac{\exp(y_i+c)}{\sum_j\exp(y_j+c)}
    =
    \frac{\exp(c)\exp(y_i)}{\exp(c)\sum_j\exp(y_j)}
    =
    \frac{\exp(y_i)}{\sum_j\exp(y_j)}.
\]

因此 $(3,1,-3)$ 与 $(103,101,97)$ 的 softmax 输出完全相同。工程实现中常利用这个性质做数值稳定化：先减去最大 logit $m=\max_j y_j$，再计算

\[
    y'_i=\frac{\exp(y_i-m)}{\sum_j\exp(y_j-m)}.
\]

这样可以避免 $\exp(y_i)$ 因为 logit 太大而溢出。这一点视频里没有展开，但它解释了为什么深度学习框架常把 softmax 与 cross-entropy 合并实现：不仅公式更简洁，数值上也更稳。

\begin{figure}[H]
    \centering
    \includegraphics[width=0.86\textwidth]{figures/fig04_softmax_example.jpg}
    \caption{softmax 的数值例子：若 logits 为 $(3,1,-3)$，取 exponential 后约为 $(20,2.7,0.05)$，归一化后得到 $(0.88,0.12,\approx 0)$。softmax 不只把值压到 $0$ 到 $1$，也会放大大 logit 与小 logit 的差距。\protect\footnotemark}
    \label{fig:softmax_example}
\end{figure}
\footnotetext{视频画面与字幕时间区间：09:26--10:08。}

\begin{importantbox}{二分类里的 sigmoid 与 softmax}
    如果只有两个类别，可以直接用二类 softmax；实际工程中更常见的是用 sigmoid 表示二分类。两者并不是互相竞争的不同哲学，而是在二分类设定下可以互相转写的等价形式。多分类时，softmax 是更自然的统一写法。
\end{importantbox}

\begin{knowledgebox}{logit 的角色：先比较，再归一化}
    softmax 的输入 $y_i$ 常被称作 logit。logit 可以很大、很小、为正、为负，它本身不需要像概率一样被限制在 $[0,1]$。softmax 先把每个 logit 指数化，再用总和归一化，因此真正影响类别结果的是\textbf{不同 logits 之间的相对大小}。如果某个 logit 比其他 logit 大很多，softmax 会把它推到接近 $1$，其他类别则接近 $0$。
\end{knowledgebox}

\begin{warningbox}{不要把 softmax 当成``让模型正确''的魔法}
    Softmax 只负责把 logits 转成一个归一化分布；它不会凭空让模型学会分类。若 logits 本来就把错误类别打得很高，softmax 只会把这个错误变成``很有信心的错误''。真正让 logits 往正确方向移动的，是后面的 loss function 与反向传播。
\end{warningbox}

% =====================================================================
\section{分类损失：MSE 与 Cross-entropy}
% =====================================================================

得到 $\boldsymbol{y}'$ 之后，就要定义它和正确标签 $\hat{\boldsymbol{y}}$ 的距离。一个样本的误差记为 $e_n$，整体 loss 是所有样本误差的平均：

\[
    L = \frac{1}{N}\sum_n e_n.
\]

问题在于：$e_n$ 应该怎么定义？本讲比较了两个候选：

\[
    e_{\text{MSE}} = \sum_i(\hat{y}_i-y'_i)^2,
\]

\[
    e_{\text{CE}} = -\sum_i \hat{y}_i\ln y'_i.
\]

\begin{figure}[H]
    \centering
    \includegraphics[width=0.86\textwidth]{figures/fig05_classification_loss.jpg}
    \caption{分类损失的两个候选：Mean Square Error 直接计算 one-hot 标签与 softmax 输出的平方距离；Cross-entropy 使用 $-\sum_i \hat{y}_i\ln y'_i$。在分类任务中，cross-entropy 是更常用的选择；从传统统计观点看，最小化 cross-entropy 等价于最大化 likelihood。\protect\footnotemark}
    \label{fig:classification_loss}
\end{figure}
\footnotetext{视频画面与字幕时间区间：12:06--13:38。}

若 $\hat{\boldsymbol{y}}$ 是 one-hot，假设正确类别为 $k$，则 cross-entropy 会简化为

\[
    e_{\text{CE}} = -\ln y'_k.
\]

也就是说，它只关心模型分给正确类别的 softmax 输出 $y'_k$：若 $y'_k$ 越接近 $1$，loss 越小；若 $y'_k$ 越接近 $0$，loss 会急剧变大。

\subsection{从 likelihood 看 cross-entropy}

如果把 softmax 输出 $y'_i$ 当作模型给第 $i$ 类的预测概率，那么一个正确类别为 $k$ 的样本，其 likelihood 就是模型分给正确类别的概率：

\[
    P(\text{correct class}=k\mid \boldsymbol{x})=y'_k.
\]

最大化 likelihood 等价于最大化 $\ln y'_k$，也等价于最小化 $-\ln y'_k$。而 one-hot 形式的 cross-entropy

\[
    -\sum_i \hat{y}_i\ln y'_i
\]

在 $\hat{y}_k=1$、其他维度为 $0$ 时正好退化成 $-\ln y'_k$。这就是老师在投影片中说的：\textbf{minimizing cross-entropy is equivalent to maximizing likelihood}。它不是``两个很像的目标''，而是同一个目标的两种写法。

\begin{knowledgebox}{为什么 PyTorch 里常常看不到 softmax}
    在 PyTorch 中，\texttt{CrossEntropyLoss} 通常把 log-softmax 与 negative log-likelihood 合在一起实现。因此模型最后一层一般直接输出 logits，不需要手动加 softmax；如果先在网络里加 softmax，再把结果送进 \texttt{CrossEntropyLoss}，就相当于把 softmax 的角色重复处理了一次。
\end{knowledgebox}

更具体地说，PyTorch 的常见写法是：

\begin{lstlisting}[language=Python]
logits = model(x)              # shape: [batch_size, num_classes]
loss = criterion(logits, y)    # y: class index, shape [batch_size]
\end{lstlisting}

其中 \texttt{y} 通常不是 one-hot 向量，而是类别编号，例如 \texttt{0, 1, 2}。这看起来和前面说的 one-hot 不同，其实只是 API 选择：框架内部知道正确类别是哪一维，就能等价地计算 $-\ln p_k$，没有必要真的在内存里展开成 one-hot。

\begin{warningbox}{一个常见实作错误：重复 softmax}
    如果模型最后一层已经写了 \texttt{softmax}，又使用 \texttt{CrossEntropyLoss}，就容易把 logits 先压成概率，再交给一个期待 logits 的 loss。这样不仅语义不清，也可能让梯度与数值稳定性变差。多分类训练时，最稳妥的习惯是：\textbf{模型输出 logits，loss 使用 cross-entropy；需要显示概率时，再在推理或评估阶段额外 softmax。}
\end{warningbox}

\begin{importantbox}{MSE 和 Cross-entropy 的共同目标}
    如果训练成功，MSE 和 cross-entropy 都希望 $\boldsymbol{y}'$ 接近 $\hat{\boldsymbol{y}}$。差别不在最终想要的答案，而在\textbf{通往答案的 error surface 长什么样}。这正是本讲标题的含义：loss function 的定义会改变 optimization 的难度。
\end{importantbox}

% =====================================================================
\section{为什么 Cross-entropy 更适合分类优化？}
% =====================================================================

老师用一个三分类例子说明 MSE 与 cross-entropy 的差异。设正确标签为

\[
    \hat{\boldsymbol{y}}=\begin{bmatrix}1\\0\\0\end{bmatrix},
\]

网络输出三个 logits $y_1,y_2,y_3$，通过 softmax 得到 $y'_1,y'_2,y'_3$。为了画出二维图，固定

\[
    y_3=-1000,
\]

只让 $y_1$ 和 $y_2$ 在 $[-10,10]$ 中变化。因为 $y_3$ 极小，softmax 后 $y'_3\approx 0$，因此主要看 $y_1,y_2$ 如何影响 loss。

当 $y_1$ 很大、$y_2$ 很小时，正确类别 class 1 的 softmax 输出 $y'_1$ 接近 $1$，loss 很小；当 $y_1$ 很小、$y_2$ 很大时，模型把错误类别看得很有信心，loss 很大。两种 loss 都满足这个大方向，但它们的坡度完全不同。

\begin{figure}[H]
    \centering
    \includegraphics[width=0.9\textwidth]{figures/fig06_loss_surface_comparison.jpg}
    \caption{MSE 与 cross-entropy 的 error surface 对比。红色区域表示大 loss，蓝色区域表示小 loss。若初始化在左上角这种``错得很有信心''的位置，MSE 的大 loss 区域反而非常平坦，gradient 很小，容易 stuck；cross-entropy 在同一区域仍有明显坡度，可以沿 gradient 往右下角的小 loss 区移动。\protect\footnotemark}
    \label{fig:loss_surface_comparison}
\end{figure}
\footnotetext{视频画面与字幕时间区间：18:56--19:23。}

\begin{warningbox}{MSE 的危险：错得很离谱时，gradient 反而可能很小}
    MSE 接在 softmax 后面时，若模型在错误类别上已经非常有信心，softmax 容易进入饱和区域。此时输出看起来离目标很远，loss 很大，但对 logits 的梯度可能很小，导致训练起步困难。也就是说，模型不是因为答案接近正确而不动，而是因为错误的 loss surface 让它在糟糕区域也走不动。
\end{warningbox}

\begin{importantbox}{Cross-entropy 的优势}
    Cross-entropy 对``把正确类别预测成接近 0''这件事惩罚非常强：$-\ln y'_k$ 会在 $y'_k\to 0$ 时迅速变大。它因此能在错得很离谱的区域保留有用的梯度，帮助模型从错误且自信的初始状态中离开。
\end{importantbox}

\subsection{把梯度写出来：差别为什么会这么大？}

上面的图给的是直觉；再把关键梯度写出来，就能看见为什么 cross-entropy 更顺。记 softmax 输出为

\[
    p_i = y'_i = \frac{\exp(y_i)}{\sum_j \exp(y_j)},\qquad
    q_i=\hat{y}_i.
\]

其中 $\boldsymbol{p}$ 是模型预测分布，$\boldsymbol{q}$ 是 one-hot 标签。对于 \textbf{softmax + cross-entropy}，

\[
    e_{\text{CE}}=-\sum_i q_i\ln p_i,
\]

它对 logit $y_i$ 的梯度会漂亮地化简成

\[
    \frac{\partial e_{\text{CE}}}{\partial y_i}=p_i-q_i.
\]

这个化简可以用一行链式法则看出来。Softmax 的导数为

\[
    \frac{\partial p_i}{\partial y_\ell}
    = p_i(\delta_{i\ell}-p_\ell),
\]

于是

\begin{align*}
    \frac{\partial e_{\text{CE}}}{\partial y_\ell}
    &= \sum_i \frac{\partial e_{\text{CE}}}{\partial p_i}
       \frac{\partial p_i}{\partial y_\ell} \\
    &= \sum_i \left(-\frac{q_i}{p_i}\right)
       p_i(\delta_{i\ell}-p_\ell) \\
    &= -q_\ell + p_\ell\sum_i q_i \\
    &= p_\ell-q_\ell,
\end{align*}

其中最后一步用到 $\boldsymbol{q}$ 是 one-hot，因此 $\sum_i q_i=1$。这个结果之所以漂亮，是因为 cross-entropy 里的 $\ln p_i$ 和 softmax 的导数正好互相配合，把复杂的 Jacobian 化成了``预测减标签''。

这条式子信息量很大：如果正确类 $k$ 的预测 $p_k$ 很小，那么 $p_k-q_k\approx -1$，梯度仍然很强；如果某个错误类 $r$ 被预测得很大，$p_r-q_r\approx p_r$，梯度会明确要求把它压下去。也就是说，\textbf{模型错得越自信，cross-entropy 给出的纠偏信号越清楚}。

MSE 接在 softmax 后面时，情况就不一样。MSE 是

\[
    e_{\text{MSE}}=\sum_i(q_i-p_i)^2.
\]

对 logit 求导时，必须再乘上 softmax 的导数：

\[
    \frac{\partial p_i}{\partial y_\ell}
    = p_i(\delta_{i\ell}-p_\ell).
\]

因此梯度中会出现 $p_i(1-p_i)$ 或 $p_ip_\ell$ 这样的因子。当 softmax 已经饱和时，某些 $p_i$ 接近 $0$ 或 $1$，这些因子就会非常小。于是会出现很尴尬的局面：\textbf{预测明明错得很严重，MSE 本身也很大，但传回 logits 的 gradient 却被 softmax 饱和压扁了}。这就是图中 MSE 左上角大 loss 区域却几乎平坦的原因。

举一个极端但很有说明力的例子。假设正确类别是 class 1，但模型 logits 给出

\[
    \boldsymbol{y}=(-10,\,10,\,-1000).
\]

Softmax 后几乎是

\[
    \boldsymbol{p}\approx(0,\,1,\,0),\qquad
    \boldsymbol{q}=(1,\,0,\,0).
\]

这代表模型非常自信地选错了 class 2。对 cross-entropy 而言，正确类梯度约为 $p_1-q_1\approx -1$，错误高分类梯度约为 $p_2-q_2\approx 1$，所以优化器会明确地提高 $y_1$、降低 $y_2$。但对 MSE 而言，softmax 已经在 $(0,1,0)$ 附近饱和，softmax Jacobian 很小，梯度信号会被削弱；这就是``错得很严重，却不好改''的核心。

\begin{importantbox}{最关键的公式对照}
    对分类任务而言，\textbf{softmax + cross-entropy} 的梯度核心是
    \[
        \frac{\partial e}{\partial y_i}=p_i-q_i,
    \]
    它直接比较``预测分布''与``正确 one-hot 标签''。而 \textbf{softmax + MSE} 的梯度还要额外乘上 softmax Jacobian，容易在饱和区变小。前者给 optimizer 的路标清楚，后者可能把路标藏在平坦的大 loss 区里。
\end{importantbox}

这不是说 MSE 在数学上``不能''用于分类；在更强的 optimizer、良好的初始化或特定任务中，它未必完全失败。老师也提到，若使用 Adam 这类自适应优化器，gradient 小时学习率可能自动被调大，训练仍有机会继续推进。但这会让训练更困难、更慢，尤其在初期 logits 已经把类别分错且 softmax 饱和时。实际分类任务中，\textbf{softmax + cross-entropy} 成为标准搭配，正是因为它给 optimization 更友善的地形。

% =====================================================================
\section{本讲综合：loss function 也是优化设计的一部分}
% =====================================================================

这一讲的核心不是背一个公式，而是建立一个更工程化的判断：\textbf{模型训练不起来，不一定只怪模型结构、learning rate 或 optimizer；loss function 本身也可能把 error surface 变得难走。}

分类任务的一条常用路线可以总结为：

\begin{enumerate}
    \item 用 one-hot vector 表示类别，避免给无序类别强加数轴距离。
    \item 让网络最后一层输出与类别数相同的 logits。
    \item 用 softmax 把 logits 变成归一化输出。
    \item 用 cross-entropy 衡量输出与 one-hot 标签的差距。
    \item 在 PyTorch 等框架中，通常让模型输出 logits，直接交给内建的 cross-entropy loss。
\end{enumerate}

从优化角度看，MSE 与 cross-entropy 的差别在于：它们虽然可能有相同的理想终点，却提供了不同的训练地形。Cross-entropy 在模型错得很有信心时仍然给出清楚的下降方向；MSE 则可能在同样区域变得平坦，使 gradient descent 一开始就迈不开步。

\begin{knowledgebox}{一句话总结}
    \textbf{改变 loss function 可以改变 optimization 的难度。}分类任务选择 cross-entropy，不只是因为它有 likelihood 的解释，更因为它让 gradient descent 面对一个更容易走的 error surface。
\end{knowledgebox}

\begin{importantbox}{和前几讲连起来看}
    前几讲讨论的是：同一个 loss surface 上，怎样靠 momentum、adaptive learning rate、learning rate scheduling 走得更好。本讲补上另一块拼图：\textbf{loss function 本身会决定 surface 长什么样}。所以训练卡住时，除了调 optimizer，也要问一个更上游的问题：我定义的目标函数，是否让正确方向在梯度上足够可见？
\end{importantbox}

\end{document}
